\chapter{Étude Technique et Choix Technologiques}
\label{ch:etude_technique}

% ------------------------------------------------------------------------------
\section*{Introduction}
\addcontentsline{toc}{section}{Introduction}
% ------------------------------------------------------------------------------
La conception d'un système d'aide à la décision pour la planification hospitalière assemble plusieurs briques : modèles de prévision, méthodes d'optimisation, service web et interface de visualisation. Ce chapitre présente les critères de sélection, puis les quatre méthodes de planification et les cinq modèles de prévision comparés. Il décrit ensuite chaque outil retenu et la raison de son choix au regard des besoins identifiés au chapitre~\ref{ch:contexte_analyse_besoins}, et les résume dans un tableau de synthèse. Les solutions non retenues sont citées comme alternatives ; elles n'ont pas fait l'objet de mesures comparatives dans le projet.

% ------------------------------------------------------------------------------
\section{Critères de Sélection}
% ------------------------------------------------------------------------------

Les technologies ont été choisies en fonction de cinq critères, déduits des besoins du projet :
\begin{itemize}
    \item \textbf{Adéquation au problème :} prévision de séries journalières, optimisation combinatoire sous contraintes et restitution interactive ;
    \item \textbf{Exécution locale :} fonctionnement sur un poste de travail, sans serveur dédié ni base de données ;
    \item \textbf{Facilité d'intégration :} bibliothèques documentées, utilisables ensemble dans un même dépôt ;
    \item \textbf{Gratuité :} outils utilisables sans licence à acquérir ;
    \item \textbf{Contrôle :} possibilité de valider les données échangées et de tester automatiquement le calcul et l'interface.
\end{itemize}

% ------------------------------------------------------------------------------
\section{Méthodes et Modèles Comparés}
% ------------------------------------------------------------------------------

\subsection{Méthodes d'Optimisation}
Le projet compare quatre méthodes sur les mêmes instances : deux heuristiques rapides et deux solveurs. Cette comparaison permet de mesurer ce qu'apporte un solveur par rapport à une règle simple.
\begin{itemize}
    \item \textbf{Stratégie naïve} : point de comparaison bas. Chaque jour, un taux d'acceptation unique s'applique à tous les services ; le personnel est partagé proportionnellement aux cibles, sans service prioritaire. Elle est déterministe et rapide, mais ne redistribue pas les postes qu'un service ne peut pas pourvoir ;
    \item \textbf{Heuristique gloutonne (\emph{Greedy})} : référence de la comparaison. Elle accepte autant d'admissions que de lits libres, puis affecte le personnel jour par jour, dans l'ordre $\text{REA} \succ \text{URG} \succ \text{CHI} \succ \text{MED}$, en deux passes : dotation minimale, puis cible. Elle est déterministe et calculée en moins de deux secondes, mais ses choix ne sont jamais remis en cause ;
    \item \textbf{Programmation linéaire en nombres entiers (MILP)} : le problème est décrit par des variables entières et binaires et par des contraintes linéaires~\cite{wolsey2020}. Le solveur fournit un \emph{gap}, écart entre la solution trouvée et une borne sur l'optimum ;
    \item \textbf{Programmation par contraintes (CP-SAT)} : elle propage des règles logiques sur des variables entières~\cite{rossi2006}. Avec plusieurs threads et un budget de temps, deux exécutions peuvent donner deux plans différents.
\end{itemize}

\begin{notebox}{« Solveur exact » ne signifie pas « optimum trouvé »}
CBC et CP-SAT sont capables de prouver l'optimalité, mais ils sont arrêtés après \TempsSolveur{} secondes. Sauf preuve terminée dans ce budget, leur résultat est une solution réalisable accompagnée d'un gap. Le statut \texttt{OPTIMAL} n'est accordé que si la recherche est allée à son terme.
\end{notebox}

\subsection{Modèles de Prévision}
Les admissions journalières de chaque service sont prévues de $J+1$ à $J+7$. Cinq modèles sont comparés~\cite{hyndman2021} :
\begin{enumerate}
    \item \textbf{Modèle naïf} : répète la dernière valeur observée ;
    \item \textbf{Modèle naïf saisonnier} : répète la valeur du même jour de la semaine précédente (période 7) ;
    \item \textbf{Moyenne mobile} : moyenne des 7 derniers jours observés ;
    \item \textbf{Lissage exponentiel simple} : moyenne à poids décroissants dans le temps ;
    \item \textbf{Régression à retards} : régression de type Ridge à coefficients positifs, utilisant les valeurs à 1, 7 et 14 jours, les moyennes des 7 et 28 derniers jours, le jour de la semaine et les jours fériés synthétiques. Les événements exceptionnels sont exclus, car ils ne seraient pas connus au moment de prévoir.
\end{enumerate}

La \textbf{régression à retards} est retenue pour les quatre services, car elle a la plus faible erreur WAPE de validation (tableau~\ref{tab:previsions_validation}). Le chapitre~\ref{ch:tests_validation_resultats} montre que cet avantage n'est pas confirmé partout sur le test final.

% ------------------------------------------------------------------------------
\section{Étude et Choix des Outils et Technologies}
% ------------------------------------------------------------------------------

% ------------------------------------------------------------------------------
\subsection{Langage de Programmation : Python}
% ------------------------------------------------------------------------------

\textbf{Python} est un langage polyvalent, largement utilisé en traitement des données et en optimisation. Sa syntaxe claire et ses nombreuses bibliothèques facilitent le développement d'une chaîne de calcul complète.

\begin{figure}[H]
    \centering
    \includegraphics[width=0.2\textwidth]{logos/python.png}
    \caption{Logo de Python.}
    \label{fig:logo_python}
\end{figure}

\paragraph{Justification du choix.} Python a été retenu pour toute la partie calcul et pour l'API, car :
\begin{itemize}
    \item il réunit dans un même langage la préparation des données, la prévision, les deux solveurs et le service web ;
    \item les calculs lourds sont confiés à des solveurs compilés, ce qui limite l'intérêt d'un langage compilé pour le reste du code ;
    \item les annotations de type sont exploitées par Pydantic et FastAPI pour contrôler les données.
\end{itemize}

% ------------------------------------------------------------------------------
\subsection{Manipulation des Données : pandas}
% ------------------------------------------------------------------------------

\textbf{pandas} est une bibliothèque Python de manipulation de données tabulaires : lecture et écriture de fichiers CSV, jointures, agrégations et gestion des dates.

\begin{figure}[H]
    \centering
    \includegraphics[width=0.2\textwidth]{logos/pandas.png}
    \caption{Logo de pandas.}
    \label{fig:logo_pandas}
\end{figure}

\paragraph{Justification du choix.} Les données du projet sont des tables datées par service, par jour et par poste, dont la plus grande compte \LignesDispo{} lignes. pandas permet de les contrôler et de les joindre sans base de données, et de conserver des résultats dans des fichiers CSV lisibles. Le validateur l'utilise aussi pour recalculer les règles à partir des tables produites.

% ------------------------------------------------------------------------------
\subsection{Apprentissage Automatique : scikit-learn}
% ------------------------------------------------------------------------------

\textbf{scikit-learn} est une bibliothèque Python d'apprentissage automatique, qui fournit des modèles de régression et de classification sous une interface commune.

\begin{figure}[H]
    \centering
    \includegraphics[height=2.5cm]{logos/scikit_learn.png}
    \caption{Logo de scikit-learn.}
    \label{fig:logo_scikit_learn}
\end{figure}

\paragraph{Justification du choix.} La régression à retards utilise sa classe \texttt{Ridge}, dont la régularisation convient à des variables corrélées comme les retards et les moyennes mobiles, et dont l'option \texttt{positive} interdit les coefficients négatifs. Le modèle reste rapide et explicable ; un réseau de neurones n'a pas été jugé nécessaire pour ces séries courtes.

% ------------------------------------------------------------------------------
\subsection{Modèles Statistiques : statsmodels}
% ------------------------------------------------------------------------------

\textbf{statsmodels} est une bibliothèque Python de modélisation statistique, qui comprend des modèles de séries temporelles.

\begin{figure}[H]
    \centering
    \includegraphics[width=0.2\textwidth]{logos/statsmodels.png}
    \caption{Logo de statsmodels.}
    \label{fig:logo_statsmodels}
\end{figure}

\paragraph{Justification du choix.} Elle fournit le lissage exponentiel simple (\texttt{SimpleExpSmoothing}), utilisé comme modèle statistique de comparaison. S'appuyer sur une implémentation existante évite de recoder l'estimation du paramètre de lissage.

% ------------------------------------------------------------------------------
\subsection{Programmation Linéaire en Nombres Entiers : PuLP et COIN-OR CBC}
% ------------------------------------------------------------------------------

\textbf{PuLP} est une bibliothèque Python de modélisation de problèmes linéaires~\cite{mitchell2011}. \textbf{CBC} est le solveur libre en nombres entiers auquel elle transmet le modèle~\cite{cbc}. Tous deux sont hébergés par le projet COIN-OR.

\begin{figure}[H]
    \centering
    \includegraphics[height=2.5cm]{logos/coin_or.png}
    \caption{Logo de COIN-OR, projet qui héberge PuLP et CBC.}
    \label{fig:logo_coin_or}
\end{figure}

\paragraph{Justification du choix.} PuLP permet d'écrire les contraintes sous une forme proche des équations du chapitre~\ref{ch:conception_systeme}, et l'installation fournit directement CBC. Le solveur donne une solution et une borne, donc un gap, sans licence commerciale. Ses limites sur ce problème, notamment avec la priorité personnel, sont mesurées au chapitre~\ref{ch:tests_validation_resultats}.

% ------------------------------------------------------------------------------
\subsection{Programmation par Contraintes : Google OR-Tools}
% ------------------------------------------------------------------------------

\textbf{Google OR-Tools} est une suite libre d'outils d'optimisation. Son solveur \textbf{CP-SAT} traite des problèmes à variables entières et binaires~\cite{ortools}.

\begin{figure}[H]
    \centering
    \includegraphics[height=2.5cm]{logos/ortools.png}
    \caption{Logo de Google OR-Tools.}
    \label{fig:logo_ortools}
\end{figure}

\paragraph{Justification du choix.} Le modèle est dominé par des variables binaires d'affectation et par des incompatibilités entre postes, auxquelles CP-SAT est adapté. Il exploite plusieurs threads (\ThreadsSolveur{} dans le projet), accepte le plan Greedy comme indication de départ et fournit une borne inférieure. Comparer CBC et CP-SAT sur la même formulation permet en outre de contrôler l'un par l'autre.

% ------------------------------------------------------------------------------
\subsection{API REST : FastAPI}
% ------------------------------------------------------------------------------

\textbf{FastAPI} est un cadre Python de développement d'API fondé sur les annotations de type~\cite{fastapi}. Il est exécuté ici par le serveur \textbf{Uvicorn}.

\begin{figure}[H]
    \centering
    \includegraphics[width=0.2\textwidth]{logos/fastapi.png}
    \caption{Logo de FastAPI.}
    \label{fig:logo_fastapi}
\end{figure}

\paragraph{Justification du choix.} Django est complet mais orienté vers des applications à base de données, absentes du projet ; Flask est léger, mais la validation des échanges y demande des extensions. FastAPI a été retenu car :
\begin{itemize}
    \item les requêtes et les réponses sont contrôlées par des schémas partagés avec le moteur de calcul ;
    \item un calcul long est soumis comme un \emph{job} et répond immédiatement ; un \texttt{ThreadPoolExecutor} à un seul travailleur évite que deux solveurs ne se disputent le processeur ;
    \item la documentation interactive (\texttt{/docs}) est générée depuis la description OpenAPI.
\end{itemize}

% ------------------------------------------------------------------------------
\subsection{Validation des Données : Pydantic}
% ------------------------------------------------------------------------------

\textbf{Pydantic} est une bibliothèque Python qui décrit des structures de données et vérifie leur conformité à l'aide des annotations de type~\cite{pydantic}.

\begin{figure}[H]
    \centering
    \includegraphics[height=2.5cm]{logos/pydantic.png}
    \caption{Logo de Pydantic.}
    \label{fig:logo_pydantic}
\end{figure}

\paragraph{Justification du choix.} Les mêmes classes (\texttt{contracts.py}) décrivent la configuration de l'hôpital, les profils, les résultats et les requêtes de l'API. Une valeur hors bornes, comme un taux d'absence supérieur à 50\,\% dans le simulateur, ou un champ inconnu est refusé avec un message qui désigne le champ fautif, sans contrôle écrit à la main.

% ------------------------------------------------------------------------------
\subsection{Structure de l'Interface Web : Next.js}
% ------------------------------------------------------------------------------

\textbf{Next.js} est un cadre de développement web fondé sur React, qui fournit la structure du projet, le serveur de développement et les outils de compilation~\cite{nextjs}.

\begin{figure}[H]
    \centering
    \includegraphics[width=0.2\textwidth]{logos/nextjs.png}
    \caption{Logo de Next.js.}
    \label{fig:logo_nextjs}
\end{figure}

\paragraph{Justification du choix.} Le tableau de bord ne se limite pas à afficher des courbes : il lance des simulations, suit leur avancement et croise plusieurs filtres. Ces interactions justifient une application dédiée. Next.js apporte une organisation de projet et un outillage TypeScript prêts à l'emploi ; la mise en page utilise \textbf{Tailwind CSS}~\cite{tailwindcss} et une feuille de style propre au projet.

% ------------------------------------------------------------------------------
\subsection{Composants d'Interface : React}
% ------------------------------------------------------------------------------

\textbf{React} est une bibliothèque JavaScript qui permet de construire une interface à partir de composants réutilisables~\cite{react}.

\begin{figure}[H]
    \centering
    \includegraphics[height=2.5cm]{logos/react.png}
    \caption{Logo de React.}
    \label{fig:logo_react}
\end{figure}

\paragraph{Justification du choix.} Les dix vues partagent les mêmes composants et les mêmes filtres. Les courbes, les barres et les cartes de chaleur sont écrites directement en composants React produisant du SVG ou des tableaux, sans bibliothèque graphique, ce qui permet de contrôler l'affichage des valeurs absentes et des plans irréalisables.

% ------------------------------------------------------------------------------
\subsection{Typage de l'Interface : TypeScript}
% ------------------------------------------------------------------------------

\textbf{TypeScript} ajoute un typage statique à JavaScript.

\begin{figure}[H]
    \centering
    \includegraphics[height=2.5cm]{logos/typescript.png}
    \caption{Logo de TypeScript.}
    \label{fig:logo_typescript}
\end{figure}

\paragraph{Justification du choix.} Les réponses de l'API sont décrites dans \texttt{api-types.ts} et utilisées par toutes les vues. Une incohérence de format, ou l'oubli d'une valeur possiblement absente, est signalée dès la vérification des types, avant l'exécution.

% ------------------------------------------------------------------------------
\subsection{Tests du Calcul et de l'API : pytest}
% ------------------------------------------------------------------------------

\textbf{pytest} est un outil de test pour Python.

\begin{figure}[H]
    \centering
    \includegraphics[height=2.5cm]{logos/pytest.png}
    \caption{Logo de pytest.}
    \label{fig:logo_pytest}
\end{figure}

\paragraph{Justification du choix.} Les \NbTestsPython{} tests couvrent les données, les prévisions, les méthodes, le validateur, les scénarios et l'API. Les \emph{fixtures} de pytest permettent de partager la configuration de référence entre les tests, et les solveurs sont vérifiés sur de petites instances. L'analyse statique est assurée par \textbf{ruff}.

% ------------------------------------------------------------------------------
\subsection{Tests de l'Interface : Playwright}
% ------------------------------------------------------------------------------

\textbf{Playwright} est un outil qui pilote un navigateur pour tester une application web.

\begin{figure}[H]
    \centering
    \includegraphics[height=2.5cm]{logos/playwright.png}
    \caption{Logo de Playwright.}
    \label{fig:logo_playwright}
\end{figure}

\paragraph{Justification du choix.} Les \NbTestsNavigateur{} parcours vérifient l'ensemble, du navigateur à l'API : ouverture des dix vues sans calcul automatique, simulation complète, cas irréalisable et changement de profil. Ils s'exécutent sur une copie des résultats. Playwright a aussi servi à produire les captures de ce rapport.

% ------------------------------------------------------------------------------
\subsection{Gestion de Versions : Git et GitHub}
% ------------------------------------------------------------------------------

\textbf{Git} est un système de gestion de versions. \textbf{GitHub} héberge le dépôt du projet.

\begin{figure}[H]
    \centering
    \includegraphics[width=0.2\textwidth]{logos/git.png}
    \hspace{1.5cm}
    \includegraphics[height=2.5cm]{logos/github.png}
    \caption{Logos de Git et de GitHub.}
    \label{fig:logo_git_github}
\end{figure}

\paragraph{Justification du choix.} Le code, les configurations, les données synthétiques, les résultats de référence et ce rapport sont versionnés ensemble : un résultat peut ainsi être rattaché à l'état du code qui l'a produit. Les simulations personnalisées, propres à chaque poste, sont exclues du dépôt.

% ------------------------------------------------------------------------------
\section{Synthèse des Choix Technologiques}
% ------------------------------------------------------------------------------

Le tableau~\ref{tab:technologies} récapitule, pour chaque fonction du système, la technologie retenue et la raison principale de son choix. Les versions sont celles installées lors de la rédaction.

\begin{table}[H]
\centering
\small
\renewcommand{\arraystretch}{1.25}
\begin{tabularx}{\linewidth}{>{\raggedright\arraybackslash}p{3.2cm} >{\raggedright\arraybackslash}p{4.6cm} >{\raggedright\arraybackslash}X}
\toprule
\textbf{Fonction} & \textbf{Technologie retenue} & \textbf{Raison principale du choix} \\
\midrule
Calcul et API & Python 3.13.2 & Un seul langage pour les données, la prévision, l'optimisation et le service \\
Données tabulaires & pandas 2.3.3 & Tables datées, fichiers CSV lisibles, sans base de données \\
Régression & scikit-learn 1.9.1 & Ridge à coefficients positifs, rapide et explicable \\
Lissage exponentiel & statsmodels 0.15.0 & Implémentation statistique existante \\
Solveur MILP & PuLP 4.3.1 / COIN-OR CBC & Formulation lisible, gap fourni, sans licence \\
Solveur par contraintes & Google OR-Tools CP-SAT 9.15 & Variables binaires, recherche parallèle, borne inférieure \\
API REST & FastAPI 0.165.0 / Uvicorn 0.54.0 & Contrats typés, documentation OpenAPI \\
Validation & Pydantic 2.14.4 & Schémas partagés, erreurs détaillées \\
Interface & Next.js 15.5.28 / Tailwind CSS 3.4.19 & Structure de projet et outillage prêts à l'emploi \\
Composants et graphiques & React 19.3.1 (SVG) & Composants partagés, contrôle des états vides \\
Typage & TypeScript 5.9.3 & Erreurs de format détectées avant l'exécution \\
Tests Python & pytest 9.2.0 & Contrôle du calcul, du validateur et de l'API \\
Tests navigateur & Playwright 1.64.1 & Parcours complets, du navigateur à l'API \\
Versions & Git / GitHub & Code, données et résultats versionnés ensemble \\
\bottomrule
\end{tabularx}
\caption{Synthèse et justification des technologies retenues.}
\label{tab:technologies}
\end{table}

% ------------------------------------------------------------------------------
\section*{Conclusion}
\addcontentsline{toc}{section}{Conclusion}
% ------------------------------------------------------------------------------
Les outils retenus sont gratuits et s'exécutent sur un poste de travail. Deux heuristiques déterministes donnent une réponse immédiate ; deux solveurs les complètent dans un budget de temps fixé, sans garantie d'atteindre l'optimum dans ce délai. Les prévisions reposent sur des modèles simples et explicables, et l'ensemble est servi par une API typée et un tableau de bord Next.js, sans base de données. Ces choix étant fixés, le chapitre suivant formalise le problème que les méthodes doivent résoudre et décrit la conception du système.
